\section{Elys 与 ChatGPT、Meta、豆包的竞争逻辑}

访谈里有一个重要判断：模型能力会逐渐平权，但 context 不平权。模型公司拥有强智能和入口，社交平台拥有关系链和内容行为，Elys 这样的创业公司拥有新范式和产品聚焦。谁能赢，不只看谁的模型强，而看谁能积累最有用、最可信、最能转化成真人连接的 context。

\subsection{四类玩家的路径}

\noindent\begin{longtable}{p{0.20\textwidth}p{0.35\textwidth}p{0.36\textwidth}}
\toprule
玩家 & 优势 & 短板 \\
\midrule
ChatGPT / OpenAI & 对话入口、个人记忆、强模型 & 社交不是主航道，娱乐和情感化产品心智未必强。 \\
Meta / 社交平台 & 关系链、内容、身份和分发 & 旧社交结构可能难以 AI-native 化。 \\
豆包/国内大厂 & C 端入口、内容生态、本土场景 & 是否把 context 视为战略核心仍不确定。 \\
Elys / 创业公司 & 从零按 AI 社交范式设计，聚焦真人连接 & 用户、信任、context 冷启动压力大。 \\
\bottomrule
\end{longtable}

\begin{warningbox}{模型公司不一定天然赢}
如果产品核心是智能问答，模型公司优势巨大；如果产品核心是真人关系、context 飞轮和社交心智，应用公司仍可能有窗口。关键是产品是否远离模型公司的主航道。
\end{warningbox}

\subsection{为什么 OpenClaw 与 Elys 属于同一叙事}

OpenClaw 让用户感觉“有个东西主动帮我干活”；Elys 让用户感觉“有个分身主动帮我面对世界”。二者都体现 proactive agent 的产品心智。区别在于，OpenClaw 面对工具和本地电脑，Elys 面对人和社交关系。

\begin{importantbox}{同一叙事：主动性}
AI 产品的根本变化不是界面从 GUI 变成 LUI，而是从 reactive 变成 proactive。模型开始主动探索、执行、带回结果，产品也因此从工具变成代理。
\end{importantbox}

\subsection{本章小结}

AI 社交竞争不只是模型能力竞争。Context、关系链、权限、产品趣味和真人连接率，都会决定最终格局。

